\documentclass[11pt]{article}

\usepackage[T1]{fontenc}
\usepackage{lmodern}
\usepackage[margin=1in]{geometry}
\usepackage{amsmath,amssymb,amsthm}
\usepackage{microtype}
\usepackage{graphicx}
\usepackage{booktabs}
\usepackage{array}
\usepackage{enumitem}
\usepackage{xcolor}
\usepackage{fancyhdr}
\usepackage[hidelinks]{hyperref}
\hypersetup{
  pdfauthor={Ioannis Tsiokos},
  pdftitle={Six Birds for Neutrino Mass: Lens-Swap Stability and Packaging Audits}
}
\usepackage{cleveref}
\crefname{section}{Section}{Sections}
\crefname{figure}{Figure}{Figures}
\crefname{table}{Table}{Tables}
\crefname{equation}{Eq.}{Eqs.}
\crefname{appendix}{Appendix}{Appendices}
\crefname{proposition}{Proposition}{Propositions}
\newtheorem{proposition}{Proposition}
\theoremstyle{remark}
\newtheorem*{remark}{Remark}
\usepackage[numbers,sort&compress]{natbib}
\usepackage{orcidlink}

\graphicspath{{figures/}}

\newif\ifpaperhasdoi
\paperhasdoitrue
\newcommand{\paperdoi}{10.5281/zenodo.23161613}
\newcommand{\manuscriptversion}{Version 2}
\newcommand{\manuscriptdate}{5 October 2026}
\newcommand{\manuscriptyear}{2026}
 
\newcommand{\Lens}{f}
\newcommand{\Completion}{U}
\newcommand{\Pack}{E}
\newcommand{\PackageOp}{\mathcal{E}}
\newcommand{\RouteMismatch}{\Delta_{\mathrm{route}}}
\newcommand{\Staging}{\mathrm{staging}}
\newcommand{\dchi}{\Delta \chi^2}
\newcommand{\mnu}{\sum m_\nu}

\numberwithin{equation}{section}

\widowpenalty=10000
\clubpenalty=10000
\raggedbottom
\emergencystretch=3em

\makeatletter
\newcommand{\affiliation}[1]{\gdef\@affiliation{#1}}
\newcommand{\printaffiliation}{\begin{center}\small \@affiliation\end{center}}
\makeatother

\title{Six Birds for Neutrino Mass:\texorpdfstring{\\}{ } Lens-Swap Stability and Packaging Audits}
\author{Ioannis Tsiokos\,\orcidlink{0009-0009-7659-5964}}
\affiliation{}
\date{Version 1: 2 March 2026\\[2pt]\manuscriptversion: \manuscriptdate\ifpaperhasdoi\\[4pt]{\small DOI: \href{https://doi.org/\paperdoi}{\paperdoi}}\fi}

\fancypagestyle{firstpage}{\fancyhf{}
  \fancyfoot[C]{\scriptsize Automorph Inc., Wilmington, DE, USA \quad \texttt{ioannis@automorph.io} \quad \textcopyright\ \manuscriptyear\ \quad CC-BY 4.0 \quad Preprint --- not peer reviewed.}
  \renewcommand{\headrulewidth}{0pt}
  \renewcommand{\footrulewidth}{0pt}
}

\begin{document}

\maketitle
\thispagestyle{firstpage}
\printaffiliation

\begin{abstract}
\label{sec:abstract}
When the SPT-3G 2018 temperature and polarization likelihood is replaced by the later SPT-3G D1 release, which observes essentially the same patch of sky, cosmological upper limits on the summed neutrino mass $\mnu$ have been reported to tighten \cite{arxiv2601_16277}. We ask a prior question: do the two released likelihoods make mutually consistent predictions? Following the Six Birds framework \cite{sixbirds_foundations,sixbirds_darkenergy}, we treat each release as a \emph{lens} on the same sky and measure, in each direction, how much worse one lens fits at parameters chosen by the other. We split the SPT quadratic part of this mismatch across spectra and multipoles with a signed allocation that respects the full covariance, so the pieces add up exactly to that part. At numerically optimized (but not certified-global) fits, the mismatch is large and strongly asymmetric. Testing D1 at the 2018 fit costs $\dchi\approx488$ for SPT with DESI DR2 BAO and $\approx63$ with a Planck-subset baseline added; the reverse direction costs $\approx79$ and $\approx23$. Where the mismatch sits depends on both baseline and direction: temperature at $2000\le\ell<3000$ leads both directions of the SPT+DESI audit, whereas polarization (EE) leads the D1-given-2018 direction once the Planck subset is included. Cutting D1 back to the 2018 multipole coverage, dropping the SPT optical-depth prior, or raising Boltzmann precision changes each tested total by less than 4\%. Our archived single-chain $\mnu$ runs, and a later multi-chain campaign, do not meet pre-declared convergence and precision criteria, and the likelihoods themselves proved sensitive to solver version and numerical settings. We therefore report no tightening magnitude. We also give a short, machine-checked (Lean 4) analysis of two Gaussian constraints with the physical gate $\mnu\ge0$. It shows when the posterior mode sits exactly at zero, that tension between two constraints that both prefer positive masses can never put it there, and that tightening one constraint raises the chance of a boundary mode only when the other constraint's mean is positive.
 \end{abstract}

\noindent\textbf{Keywords:} neutrino mass; cosmological bounds; SPT-3G; likelihood stability; cross-lens audit; covariance accounting; formal verification; Six Birds framework

\section{Introduction}
\label{sec:intro}

The summed neutrino mass $\mnu$ is one of the main targets of precision cosmology. Its constraints come from combining early-universe information in the cosmic microwave background (CMB) with late-time distance measurements such as baryon acoustic oscillations (BAO). As upper limits approach the minimum allowed by oscillation experiments, a practical question becomes pressing: when a limit moves, how much of the movement reflects the sky, and how much reflects choices in how the data were processed and compressed into a likelihood?

A sharp example is reported in \cite{arxiv2601_16277}. The authors compare two releases of the South Pole Telescope SPT-3G temperature and polarization (TT/TE/EE) likelihood: the 2018 release and the later D1 release built from the 2019--2020 seasons. Holding the rest of their CMB baseline fixed (a subset of Planck 2018 products plus Planck PR4 lensing, together with the matching SPT lensing product), they find that switching from the 2018 to the D1 combination tightens the $\mnu$ upper limit and moves the posterior towards the physical boundary $\mnu=0$. They call the shift ``curious'', because both releases observe essentially the same field, and they point out that the decisive test would be to reprocess the 2018 raw data through the D1 pipeline. We do not claim that test as our idea, and we do not perform it.

\paragraph{The question we ask.}
Before interpreting a parameter shift, one can ask whether the two released likelihoods agree with each other at all. We frame this with the Six Birds framework \cite{sixbirds_foundations,sixbirds_darkenergy}. Each likelihood release is a \emph{lens}: a map from the sky (together with instrument and foreground nuisances) to a likelihood object. Here ``lens'' never means gravitational lensing. A model class with its nuisance parameterization and priors is a \emph{completion}. The posterior we report is the \emph{package} obtained by composing the two. If two lenses describe the same sky consistently, parameters chosen by one should fit the other reasonably well. We measure how far this fails, in each direction separately, and where in the data the failure sits.

\paragraph{What this paper contributes.}
\begin{enumerate}[leftmargin=*, itemsep=2pt]
  \item \textbf{Directional cross-lens audits.} We fit one SPT release (with a common completion), transfer the shared cosmological parameters to the other release, and record the loss in fit relative to that release's own best fit. We do this in both directions, for two baselines (SPT with DESI DR2 BAO, and SPT with a Planck subset, Planck PR4 lensing and DESI), and under three controls: a multipole-support cut, removal of the SPT optical-depth prior, and higher Boltzmann precision (\cref{sec:results-audits}).
  \item \textbf{Exact localization with correlated data.} We split the SPT quadratic part of each mismatch over spectra and multipole bands with a signed allocation that uses the full covariance matrix. The pieces sum exactly to that quadratic part, unlike block-by-block chi-squares, which fail to add up when blocks are correlated; native prior and normalization terms are accounted for separately (\cref{sec:ledger}). The resulting maps show that the location of the mismatch depends on both the baseline and the direction of transfer.
  \item \textbf{A clear statement of what the $\mnu$ chains support.} Our sampled posteriors do not pass convergence and quantile-precision checks fixed in advance, and the likelihood evaluations turned out to be sensitive to solver version and numerical settings. We therefore report no magnitude for the $\mnu$ tightening, and we explain what a qualified claim would require (\cref{sec:results-mnu}).
  \item \textbf{A verified boundary mechanism.} For two Gaussian constraints and the gate $\mnu\ge0$, we state exactly when the posterior mode sits at zero, show that two constraints that both prefer positive masses can never put it there, and show that tightening one constraint can make a boundary mode either more or less likely depending on the other constraint's sign. All statements are checked in Lean~4 with Mathlib \cite{lean4,mathlib} (\cref{sec:boundary-mode} and Appendix~\ref{app:formal}).
\end{enumerate}

\paragraph{What this paper does not claim.}
We do not identify the instrumental or analysis cause of the disagreement, and we do not claim a correct value or bound for $\mnu$. Our audit values come from numerical optimizations that are not certified to be global optima, and they are not calibrated significance tests. The SBT language provides the vocabulary of the audits. It is not a theorem that supplies the missing physical steps.

\paragraph{Outline.}
\Cref{sec:related-work} places the work in context. \Cref{sec:sbt-framework} introduces lenses, packages and the directional audit. \Cref{sec:baselines} defines the data combinations and \cref{sec:methods} the audit and sampling procedures. \Cref{sec:results-audits} presents the audits and \cref{sec:results-mnu} the status of the $\mnu$ posteriors. \Cref{sec:boundary-mode} analyses boundary modes, and \cref{sec:discussion,sec:conclusion} discuss limitations and conclude. The appendices list the formal results, describe reproducibility, and record the version history.
 \section{Related work}
\label{sec:related-work}

The lens-swap observation that motivates this work is reported in \cite{arxiv2601_16277}, which also discusses its connection to the preference of DESI BAO for a lower matter density. The Six Birds vocabulary of lenses, completions, packaging and audits comes from \cite{sixbirds_foundations}, and its use for train/test audits of cosmological data combinations follows \cite{sixbirds_darkenergy}. Our cross-lens audits are related to posterior predictive checks \cite{gelman_meng_stern1996_ppc}; here we compare best-fit candidates rather than full predictive distributions. The computations use Cobaya \cite{cobaya} with CLASS \cite{class} or CAMB \cite{camb}, the SPT-3G likelihoods through \texttt{candl} \cite{candl,spt3g2018_ttteee,spt3g_d1}, Planck products \cite{planck2018_cosmoparams,planck2018_likelihood,planck_pr4_lensing}, and DESI DR2 BAO \cite{desi_dr2_bao_docs,bao_data_v26,desi_dr2_results_ii_bao}. Convergence diagnostics follow \cite{gelman_rubin1992,vehtari2021_rhat}, and formal proofs use Lean~4 and Mathlib \cite{lean4,mathlib}.
 \section{Lenses, packages and directional audits}
\label{sec:sbt-framework}

This section introduces the few Six Birds / SBT notions we need \cite{sixbirds_foundations,sixbirds_darkenergy} and defines the audit statistic used throughout. The framework supplies vocabulary and a checklist. The quantitative content of the paper lies in the statistics defined here and in the mathematics of \cref{sec:boundary-mode}.

\paragraph{Lenses, completions and packages.}
Let $X$ denote the underlying state: the cosmology together with instrumental, calibration and foreground degrees of freedom. A \emph{lens} is a data-processing and compression map
\begin{equation}
  \Lens : X \to \mathcal{L},
  \label{eq:lens_map}
\end{equation}
which ends in a likelihood object. In this paper a lens is a released likelihood, such as SPT-3G 2018 or SPT-3G D1. A \emph{completion} $\Completion$ supplies what the lens leaves open: here the $\Lambda$CDM$+\mnu$ model with three degenerate massive neutrinos, the physical gate $\mnu\ge0$, each release's own nuisance model and priors, and any external data held fixed (for example DESI BAO). The reported inference, a posterior or a best fit, is the \emph{package}
\begin{equation}
  \Pack \;=\; \Completion\circ\Lens .
  \label{eq:packaging_def}
\end{equation}
If two lenses view the same sky coherently, swapping one for the other should not drastically change the package. When the package does change, SBT asks two questions before any physical interpretation: do the lenses predict each other, and does the change survive controls on staging (cuts, support, resolution)?

\paragraph{The directional audit.}
Let $\log\mathcal{L}_A$ and $\log\mathcal{L}_B$ be the log-likelihoods of lenses $A$ and $B$, each including the common completion, as functions of a parameter vector $\theta$. Let $\hat\theta_A$ and $\hat\theta_B$ be best fits. The mismatch of $B$ given $A$ is
\begin{equation}
  \dchi_{B|A} \;\equiv\; -2\left[\log \mathcal{L}_B(\hat\theta_A) - \log \mathcal{L}_B(\hat\theta_B)\right].
  \label{eq:dchi_def}
\end{equation}
It measures how much worse $B$ fits at parameters chosen by $A$ than at its own best fit. Two properties matter.

\begin{itemize}[leftmargin=*, itemsep=2pt]
  \item \emph{Sign.} $\dchi_{B|A}\ge0$ holds whenever $\log\mathcal{L}_B(\hat\theta_B)\ge\log\mathcal{L}_B(\hat\theta_A)$. A global maximum over a domain that contains $\hat\theta_A$ guarantees this, but an optimizer reporting success does not. In practice we start the reference fit from the transferred point and keep the better of the two, so the ordering holds between the retained points, without any claim that the global maximum has been found. (Lean: Appendix~\ref{app:formal}, \texttt{Audit}.)
  \item \emph{Direction.} In general $\dchi_{B|A}\ne\dchi_{A|B}$. The two lenses have different noise, multipole coverage and nuisance freedom, so ``how well 2018 predicts D1'' and ``how well D1 predicts 2018'' are different questions. We always report both.
\end{itemize}

In practice the two lenses share only some parameters. We write $\theta=(\phi,\eta_B)$, where $\phi$ are the shared parameters transferred from the train fit and $\eta_B$ are nuisance parameters of the test lens that are re-fitted (profiled). The audit then compares the profiled likelihood $\max_{\eta_B}\log\mathcal{L}_B(\hat\phi_A,\eta_B)$ with the reference fit of $B$. The statistic is a comparative diagnostic in native $-2\Delta\log\mathcal{L}$ units. It is not a calibrated $p$-value: the two releases differ in seasons, noise and nuisance models, and translating $\dchi$ into a significance would require a declared null distribution, which we do not construct.

\begin{figure}[t]
  \centering
  \includegraphics[width=\linewidth]{figures/fig_audit_schema.pdf}
  \caption{The directional audit $\dchi_{B|A}$. Each lens is fitted together with the same completion. The shared parameters $\hat\phi_A$ from the train lens are transferred to the test lens, whose remaining free nuisance parameters (if any) are re-fitted. The loss in fit relative to the test lens's own best fit is recorded; its SPT quadratic part is split over spectra and multipole bands, and native prior and normalization terms are kept as separate ledger entries (\cref{sec:ledger}). Exchanging the roles of $A$ and $B$ gives $\dchi_{A|B}$.}
  \label{fig:audit-schema}
\end{figure}

\paragraph{Staging controls.}
SBT separates genuine disagreement from \emph{staging}: changes in cuts, support or resolution that alter a package even when every component is internally consistent \cite{sixbirds_foundations}. For CMB likelihoods the obvious staging variable is multipole coverage. D1 covers TT down to $\ell=400$ and TE/EE up to $\ell=4000$, which 2018 does not. A support-cut control removes this extra D1 coverage and repeats the audit (\cref{sec:common-staging}). If the mismatch survives, extra support cannot explain it. A support cut does not equalize window functions, binning or resolution, so it tests one specific staging explanation and not all of them.

\paragraph{The six primitives in this setting.}
\Cref{tab:sixbirds_primitives} lists how the six SBT primitives appear here. The table is a reading guide, not an additional result.

\begin{table}[t]
  \centering
  \caption{The six Six Birds primitives and their role in the SPT 2018 $\leftrightarrow$ D1 lens swap.}
  \label{tab:sixbirds_primitives}
  \setlength{\tabcolsep}{4pt}
  \renewcommand{\arraystretch}{1.08}
  \small
  \begin{tabular}{@{}>{\raggedright\arraybackslash}p{0.17\linewidth} >{\raggedright\arraybackslash}p{0.27\linewidth} >{\raggedright\arraybackslash}p{0.50\linewidth}@{}}
    \toprule
    Primitive & Informal meaning & Role in this work \\
    \midrule
    \textbf{P1 Rewrite} & Minimal correction that restores consistency of a package &
    Not applied to the SPT data. A toy template projection (Appendix~\ref{app:reproducibility}) reduces a fitted residual by construction; this is descriptive fit improvement, not evidence of predictive repair. \\
    \textbf{P2 Gating} & Constraints defining admissible objects &
    Physical gate $\mnu\ge0$; finite prior box $\mnu\le5$\,eV; choice of included likelihoods. \\
    \textbf{P3 Route mismatch} & Packages that depend on the route taken &
    The 2018 and D1 lenses give different packages; quantified by $\dchi_{B|A}$ and $\dchi_{A|B}$. \\
    \textbf{P4 Staging} & Dependence on cuts, support, resolution &
    Support-cut control removing D1's extra multipole coverage. \\
    \textbf{P5 Packaging} & The package $\Completion\circ\Lens$ &
    $\mnu$ posteriors and best fits, together with nuisance and native priors (including each release's own optical-depth prior). \\
    \textbf{P6 Audit} & Held-out prediction and accounting &
    Directional audits and exact signed localization by spectrum and multipole. \\
    \bottomrule
  \end{tabular}
\end{table}
 \section{Data and baselines}
\label{sec:baselines}

We aim to change one CMB lens while holding everything else as fixed as possible. Throughout, $A$ denotes the SPT-3G 2018 lens and $B$ the SPT-3G D1 lens. We write $\dchi_{\mathrm{D1}|2018}$ for $\dchi_{B|A}$ (D1 tested at the 2018 fit) and $\dchi_{2018|\mathrm{D1}}$ for the reverse.

\subsection{The two SPT-3G lenses}
\label{sec:spt-lenses}

Both lenses are multifrequency TT/TE/EE likelihoods evaluated with \texttt{candl} \cite{candl}:
\begin{itemize}[leftmargin=*, itemsep=2pt]
  \item \textbf{SPT 2018} \cite{spt3g2018_ttteee} (\texttt{candl\_data.SPT3G\_2018\_TTTEEE\_multifreq}): TT for $750\le\ell\le3000$, TE and EE for $300\le\ell\le3000$.
  \item \textbf{SPT D1} \cite{spt3g_d1} (\texttt{spt\_candl\_data.SPT3G\_D1\_TnE\_multifreq}): TT for $400\le\ell\le3000$, TE and EE for $400\le\ell\le4000$.
\end{itemize}
These ranges follow \cite{arxiv2601_16277}. Each release carries its own nuisance model and internal priors, and we keep them as part of the lens. In particular, the two releases include different Gaussian priors on the optical depth $\tau$: mean $0.054$ and standard deviation $0.0074$ for 2018, and mean $0.051$ and standard deviation $0.006$ for D1. A lens swap therefore changes this prior as well as the data. \Cref{sec:results-audits} tests how much this matters.

\subsection{Planck subset and lensing}
\label{sec:planck-subset}

To mirror the CMB baseline of \cite{arxiv2601_16277}, we add a subset of Planck 2018 (PR3) products and Planck PR4 lensing \cite{planck2018_likelihood,planck_pr4_lensing,planck_pr4_lensing_maps}:
\begin{itemize}[leftmargin=*, itemsep=2pt]
  \item low-$\ell$ TT and EE for $\ell<30$ (EE via the \texttt{EE\_sroll2} variant in Cobaya);
  \item high-$\ell$ TT restricted to $30\le\ell\le756$ (\texttt{planck\_2018\_highl\_plik.\allowbreak TT\_lite\_native} with \texttt{bins\_for\_L\_range=\allowbreak[30,756]});
  \item no Planck high-$\ell$ TE or EE;
  \item Planck PR4 lensing (\texttt{planckpr4lensing.PlanckPR4Lensing}).
\end{itemize}
This is not the full Planck 2018 likelihood \cite{planck2018_cosmoparams}. Unlike \cite{arxiv2601_16277}, we include no SPT lensing likelihood: neither the 2018 lensing spectrum nor the D1-era MUSE lensing spectrum that they swap in alongside D1. Our ``full CMB'' baseline therefore swaps only the SPT TT/TE/EE likelihood. \Cref{tab:baseline-defs} compares the two definitions.

\begin{table}[t]
  \centering
  \caption{CMB baselines as defined in \cite{arxiv2601_16277}. Our ``full CMB'' baseline uses the same SPT primary multipole ranges and Planck subset, keeps Planck PR4 lensing, and omits both SPT lensing likelihoods, so that only the SPT TT/TE/EE lens changes.}
  \label{tab:baseline-defs}
  \small
  \begin{tabular}{@{} p{0.14\linewidth} p{0.22\linewidth} p{0.31\linewidth} p{0.25\linewidth} @{}}
\toprule
Baseline & SPT primary ($\ell$ ranges) & Planck PR3 primary subset & Lensing \\
\midrule
CMB & TT $[750,3000]$; TE/EE $[300,3000]$ & TT $\ell<30$ (Commander); EE $\ell<30$ (SRoll2); TT $30<\ell<756$ (Plik) & SPT-3G 2018 lensing + Planck PR4 lensing \\
\addlinespace
CMB-D1 & TT $[400,3000]$; TE/EE $[400,4000]$ & as above & SPT-3G 2019--2020 MUSE lensing + Planck PR4 lensing \\
\bottomrule
\end{tabular}
 \end{table}

\subsection{DESI DR2 BAO}
\label{sec:desi-bao}

For BAO we use a Gaussian likelihood built from the public DESI DR2 BAO mean vector and covariance for the combined ``ALL'' sample \cite{desi_dr2_bao_docs,bao_data_v26,desi_dr2_results_ii_bao,desi_dr2_results_i_lya}. It includes $D_M/r_d$, $D_H/r_d$ and $D_V/r_d$ points from galaxy and Lyman-$\alpha$ forest BAO \cite{eisenstein2005_baopeak}. Here $D_M=(1+z)D_A$, $D_H=c/H(z)$ and $D_V=(zD_HD_M^2)^{1/3}$. Predictions use the same Boltzmann code as the CMB likelihoods in each run.

\subsection{Two audit baselines}
The audits and posteriors use two completions:
\begin{itemize}[leftmargin=*, itemsep=2pt]
  \item \textbf{SPT+DESI}: one SPT lens plus DESI DR2 BAO, computed with CAMB.
  \item \textbf{Full CMB+DESI}: one SPT lens plus the Planck subset, Planck PR4 lensing and DESI DR2 BAO, computed with CLASS.
\end{itemize}
Both use $\Lambda$CDM$+\mnu$ with three degenerate massive neutrinos, a flat prior on $\mnu\in[0,5]$\,eV, and the six standard $\Lambda$CDM parameters. The two baselines also differ in Boltzmann code and in which foreground parameters are free. A difference between them therefore cannot be attributed to the Planck subset alone.
 \section{Methods}
\label{sec:methods}

\subsection{Two audit scopes}
\label{sec:audit-metrics}

We run the directional audit of \cref{eq:dchi_def} at two levels.

\paragraph{Cosmological candidate transfer (main audit).}
Both SPT likelihoods and the completion are placed in a single Cobaya model with one shared Boltzmann calculation, so the two lenses see identical theory spectra at identical parameters. The model rejects any difference between the two sides in theory settings, sampled parameterization, hard prior bounds or completion components. Each lens is fitted by maximizing the sum of its own native log-likelihood and the completion's log-likelihoods within the hard prior bounds, without the separate Cobaya prior density. The native \texttt{candl} priors, including the optical-depth priors, remain part of each lens. The transferred vector $\phi$ contains the seven cosmological parameters, including $\mnu$. In the full CMB+DESI baseline the test lens also re-fits the Planck calibration $A_{\mathrm{planck}}$ and the tSZ and kSZ amplitudes. In the SPT+DESI baseline no parameters are re-fitted. Other SPT nuisance parameters stay at their packaged default values. The joint $\dchi$ includes the completion terms, and we report the SPT and completion contributions separately.

The fits use bounded derivative-free optimization (BOBYQA) from several starting points. The reference fit for the test lens is started from, among other points, the transferred candidate, and the better candidate is retained. All endpoints are re-evaluated with a freshly built model, which reproduces the native likelihoods, spectra and accounting to better than $10^{-10}$ in $\chi^2$. The fits are numerical candidates and are not certified global optima.

\paragraph{Fixed-spectrum nuisance transfer (secondary).}
To isolate nuisance packaging, we also run the audit at fixed theory spectra: the packaged test spectra shipped with each release, not cosmologies fitted to each survey. The tSZ and kSZ amplitudes are shared between the lenses. In addition, a capped set of twelve calibration and foreground parameters of the test lens is re-fitted. The transferred profile starts from the transferred point. The reference fit releases both the shared and the capped parameters and is searched from two starting points: the transferred profile and the lens's own fit. The capped sets are listed in Appendix~\ref{app:reproducibility}. This scope probes how far nuisance re-tuning can absorb the disagreement between the packaged spectra. It does not test cosmological cross-prediction.

\subsection{Exact localization with correlated data}
\label{sec:ledger}

For a Gaussian likelihood with residual $r=d-t(\theta)$ and covariance $C$, the quadratic form is $Q=r^\top C^{-1}r$. A tempting localization inverts each block of $C$ separately and sums the block chi-squares. When blocks are correlated this sum differs from $Q$. For example, with
\[
  C=\begin{pmatrix}1&0.9\\0.9&1\end{pmatrix},\qquad r=\begin{pmatrix}1\\2\end{pmatrix},
\]
the full form is $Q\approx7.37$, while the two one-dimensional block chi-squares sum to $1+4=5$. Instead we allocate the full form to coordinates,
\begin{equation}
  q_i \;=\; r_i\,(C^{-1}r)_i, \qquad \sum_i q_i = Q,
  \label{eq:signed_alloc}
\end{equation}
which shares each cross term $r_iC^{-1}_{ij}r_j$ equally between $i$ and $j$. Summing $q_i$ over any partition into groups $g$ (here spectrum $\times$ multipole band) gives group contributions that add up exactly to $Q$. The audit uses the difference of these allocations between the transferred and reference endpoints,
\begin{equation}
  \Delta Q_g \;=\; \sum_{i\in g}\bigl[q_i(\hat\theta_A)-q_i(\hat\theta_B)\bigr],
  \label{eq:block_dchi}
\end{equation}
where each endpoint uses its own effective covariance. This matters because the SPT 2018 beam covariance depends on the parameters, and a Hartlap correction applies where present. Individual $q_i$ and $\Delta Q_g$ can be negative. They are an accounting of one quadratic form and not independent chi-square variables. Only the full form $Q$ is guaranteed to be nonnegative.

The SPT part of the audit then reconciles exactly with the native likelihood:
\begin{equation}
  \dchi_{\mathrm{SPT}} \;=\; \sum_g\Delta Q_g \;+\; \Delta Q_{\mathrm{unlocalized}} \;+\; \Delta(\text{native prior penalty}) \;+\; \Delta(\log\det C).
  \label{eq:ledger}
\end{equation}
Bins outside the plotted bands are kept as an explicit ``unlocalized'' term. For example, the 2018 TE and EE bins with $300\le\ell<400$ fall below the first band edge. The exact-algebra parts of this ledger (solve, allocation and grouping) are verified in Lean (Appendix~\ref{app:formal}, \texttt{CovarianceLedger}). The numerical reconstruction is checked against the native likelihood in every run.

\subsection{Controls}
\label{sec:common-staging}
We repeat the cosmological audits under three controls, each of which re-runs all four fits:
\begin{itemize}[leftmargin=*, itemsep=2pt]
  \item \textbf{Support cut.} D1 is restricted to TT $750\le\ell\le3000$ and TE/EE $\ell\le3000$, removing its extra coverage, while the 2018 lens is left unchanged. D1 cannot supply the 2018 TE/EE modes with $300\le\ell<400$, and binning and window functions are not matched.
  \item \textbf{Optical-depth prior} (full CMB+DESI). The native SPT Gaussian prior on $\tau$ is removed from both lenses; all other priors and the Planck low-$\ell$ EE likelihood are kept.
  \item \textbf{Boltzmann precision} (full CMB+DESI). Seven CLASS integration, sampling and quadrature settings are raised towards the upstream reference values, with the physical model unchanged.
\end{itemize}

\subsection{Posterior sampling and diagnostics}
\label{sec:mnu-inference}

We sample the seven parameters $\omega_b$, $\omega_c$, $H_0$, $\ln(10^{10}A_s)$, $n_s$, $\tau$ and $\mnu$ with Cobaya's Metropolis sampler. The full CMB+DESI runs also sample $A_{\mathrm{planck}}$ and the tSZ and kSZ amplitudes. The SPT+DESI runs keep all SPT nuisance parameters at their packaged defaults. For CLASS we sample an internal mass parameter, identical to $\mnu$, and pass $m_{\mathrm{ncdm}}=(\mnu/3,\mnu/3,\mnu/3)$ for three degenerate species. The identity between the sampled and reported mass is pure bookkeeping. Passing the species masses with seventeen significant digits, not eight decimal places, makes the numerical mapping faithful to the intended masses.

Empirical quantiles are computed from the inverse weighted empirical CDF, with integer holding times treated exactly. Such quantiles do not depend on how identical samples are split into rows. Convergence is assessed with independent chains using rank-normalized split $\hat R$, bulk and tail effective sample sizes, and the Monte Carlo standard error (MCSE) of the 95th-percentile, against thresholds fixed in advance \cite{gelman_rubin1992,vehtari2021_rhat}. A sampled parameter passes when $\hat R\le1.01$; the bulk, tail and quantile ESS are all at least 400; the quantile MCSE is positive and within its limit; and the chronological drift of the quantile and its sensitivity to equal-length chain truncation are within fixed multiples of that limit. For $\mnu$ the MCSE limit applies to the 95th percentile and is absolute: $0.005$\,eV in the first SPT+DESI campaign and $0.001$\,eV in the full CMB+DESI campaign and in all later SPT+DESI campaigns. For other parameters it is $5\%$ of the empirical posterior standard deviation. A group of chains qualifies only if every sampled parameter passes. Chains from different temperatures, solver versions or numerical settings are never pooled. These thresholds are empirical checks, not convergence proofs. History-dependent proposal adaptation is in particular not covered by any theorem we invoke.

In the 28 default-prior chain families examined for this purpose, the largest saved post-burn-in mass is about $0.36$\,eV, far from the prior box $\mnu\le5$\,eV. (Some later trial cohorts with modified numerical settings reached the box during unconverged start-up phases.) Lack of contact with the upper gate does not show that a wider prior would leave the 95th percentile unchanged, as \cref{prop:upper-gate} makes precise.
 \section{Audit results}
\label{sec:results-audits}

\Cref{tab:audit-summary} and \cref{fig:audit-totals} collect every directional audit. \Cref{fig:localization} shows where the mismatch sits for the two cosmological baselines.

\begin{table}[t]
  \centering
  \caption{Directional mismatch for SPT 2018 ($A$) and SPT D1 ($B$). $\dchi_{\mathrm{D1}|2018}$ evaluates D1 at the 2018 fit. Cosmological rows give the joint value (SPT plus completion) at retained numerical candidates. The fixed-spectrum row is the SPT-only nuisance-transfer diagnostic of \cref{sec:audit-metrics}. In the full CMB+DESI block, the three controls were started from the warm-start baseline fits and should be compared with that row.}
  \label{tab:audit-summary}
  \small
  \begin{tabular}{@{}l l r r@{}}
\toprule
Audit scope & Variant & $\dchi_{\mathrm{D1}|2018}$ & $\dchi_{2018|\mathrm{D1}}$ \\
\midrule
SPT only, fixed spectra & nuisance-profiled & 849.63 & 414.39 \\
 & + support cut & 837.15 & 415.79 \\
\addlinespace
SPT + DESI & baseline & 488.30 & 79.13 \\
 & + support cut & 486.61 & 80.52 \\
\addlinespace
Full CMB + DESI & baseline (multistart) & 62.59 & 23.17 \\
 & baseline (warm start) & 62.39 & 23.19 \\
 & + support cut & 64.62 & 23.39 \\
 & $-$ SPT $\tau$ prior & 61.28 & 23.11 \\
 & + CLASS precision & 61.39 & 23.08 \\
\bottomrule
\end{tabular}
 \end{table}

\begin{figure}[t]
  \centering
  \includegraphics[width=0.9\linewidth]{figures/fig_audit_totals.pdf}
  \caption{Directional mismatch for every audit (values in \cref{tab:audit-summary}; logarithmic axis). Blue: D1 tested at the 2018 fit. Orange: 2018 tested at the D1 fit. In every setting the D1-given-2018 direction is the larger one, and no control moves either direction by more than a few percent.}
  \label{fig:audit-totals}
\end{figure}

\subsection{The mismatch is large and asymmetric}
With SPT and DESI alone, transferring the 2018 cosmology to D1 costs $\dchi_{\mathrm{D1}|2018}=488.3$, of which the SPT likelihood contributes $487.8$ and BAO $+0.5$. The reverse transfer costs $\dchi_{2018|\mathrm{D1}}=79.1$ (SPT $79.6$, BAO $-0.5$). With the Planck subset and PR4 lensing added, both values are smaller, $62.6$ and $23.2$. This is consistent with the additional shared data holding the two fits closer together, although the baselines also differ in Boltzmann code and free foregrounds. The asymmetry remains (SPT contributions $60.6$ and $25.1$). The fixed-spectrum audit gives the same picture in nuisance space: even after the test lens re-fits twelve calibration and foreground parameters, the mismatch is $849.6$ and $414.4$. Within that run, re-fitting lowers the transferred mismatch from $2214.8$ to $849.6$ ($62\%$) and from $813.5$ to $414.4$ ($49\%$). Nuisance freedom thus absorbs a large share of the disagreement, but a substantial remainder survives the capped re-fit.

These numbers are differences between retained numerical candidates. In one instance, a later search found a better 2018 source fit in the full CMB baseline at higher CLASS precision, and the forward mismatch at that source changed by $4.5\%$ (from $61.4$ to $64.2$). Optimizer choices therefore move the values at the level of a few percent. This is far below the asymmetry and far below the size of the effects themselves. We do not convert these values into significances.

\subsection{Where the mismatch sits depends on baseline and direction}

\begin{figure}[t]
  \centering
  \includegraphics[width=\linewidth]{figures/fig_localization.pdf}
  \caption{Signed localization of the SPT quadratic part of each audit (\cref{eq:signed_alloc,eq:block_dchi}) by spectrum and multipole band. Columns are bands starting at the indicated $\ell$, each extending to the next (the last band covers $3000\le\ell<4000$). Red cells add to the mismatch and blue cells reduce it; each panel has its own symmetric color scale. The title gives the total SPT quadratic difference $\Delta Q$. Cells sum exactly to it, apart from bins below $\ell=400$, which are reported separately ($2.0$ in the SPT+DESI 2018-given-D1 panel and below $0.05$ elsewhere). The numbers are allocations of one correlated quadratic form, not independent chi-squares.}
  \label{fig:localization}
\end{figure}

With SPT and DESI, high-multipole temperature leads in both directions (\cref{fig:localization}, top). Over $2000\le\ell<3000$, TT contributes $291.3$ of the D1-given-2018 mismatch and EE $121.5$. In the reverse direction TT contributes $74.5$ and EE $0.7$. The single largest band in both directions is TT at $2500\le\ell<3000$. The fixed-spectrum audit agrees: there the same TT band contributes $620.5$ and $375.3$ in the two directions.

With the Planck subset added, the picture changes (\cref{fig:localization}, bottom). In the D1-given-2018 direction, polarization leads. Over $2000\le\ell<3000$, EE contributes $36.1$, TT $14.1$ and TE $4.7$, and the largest bands are EE at $2000\le\ell<2500$ ($23.5$) and $2500\le\ell<3000$ ($12.6$). In the reverse direction, TT still leads over $2000\le\ell<3000$ ($12.6$, against $0.9$ for EE), with its largest band at $2500\le\ell<3000$ ($10.4$). Summed over all bands, however, the reverse direction receives comparable TT and EE contributions ($11.3$ and $11.7$), because EE adds at intermediate multipoles. This direction-dependent pattern persists under the support cut, without the $\tau$ prior, at higher CLASS precision, and at the improved 2018 source fit mentioned above.

A statement that the disagreement ``lives in high-$\ell$ TT'' is therefore true for the SPT+DESI audits but not in general. Because the two baselines also differ in Boltzmann code and in free foreground parameters, this comparison does not show that the Planck subset alone causes the change. All of these multipoles lie inside the range covered by both lenses. The allocations identify where the packaged likelihood objects disagree. They do not identify which part of the processing (calibration, beams, foreground modeling, covariance) is responsible.

\subsection{Controls change the totals by a few percent}
Cutting D1 back to the 2018 multipole coverage changes the SPT+DESI mismatch by $-0.35\%$ and $+1.8\%$ in the two directions. It changes the full CMB+DESI values by $+3.6\%$ and $+0.9\%$ relative to the warm-start baseline, and the fixed-spectrum values by $-1.5\%$ and $+0.3\%$. Removing the native SPT optical-depth priors changes the full CMB+DESI values by $-1.8\%$ and $-0.4\%$, and raising CLASS precision changes them by $-1.6\%$ and $-0.5\%$. None of these controls changes the leading localization. Extra D1 multipole coverage, the differing $\tau$ priors and default Boltzmann precision therefore do not explain the mismatch at the audited candidates. Each control is a separate numerical refit, so these small differences combine the effect of the control with optimizer variation. They should not be read as exact isolated effects.
 \section{What the sampled \texorpdfstring{$\mnu$}{neutrino-mass} posteriors show}
\label{sec:results-mnu}

The audits compare best fits. The quantity of physical interest is the marginal posterior of $\mnu$ and its 95th percentile, which is the usual upper limit. This section asks what our sampling supports about how the limit changes when D1 replaces 2018.

\subsection{Archived single-chain runs}
Our first production runs were one Metropolis chain per baseline and lens. \Cref{tab:archived-chains} lists their empirical 95th percentiles together with diagnostics computed with chronological half-splits and exact holding-time weights.

\begin{table}[t]
  \centering
  \caption{Archived single-chain $\mnu$ runs. Percentiles are empirical summaries of the stored samples after removing the first 20\% of rows. ``First half'' and ``second half'' are chronological halves of the retained chain. Split $\hat R$ and the effective-sample-size proxy are classical scalar diagnostics. None of these runs meets the convergence criteria of \cref{sec:mnu-inference}, so the percentiles are not posterior bounds.}
  \label{tab:archived-chains}
  \small
  \begin{tabular}{@{}l r r r r r@{}}
\toprule
Archived single chain & \multicolumn{3}{c}{empirical 95th percentile [eV]} & split $\hat R$ & ESS proxy \\
\cmidrule(lr){2-4}
 & all & first half & second half & & \\
\midrule
SPT 2018 + DESI & 0.195 & 0.084 & 0.217 & 1.86 & 4.8 \\
SPT D1 + DESI & 0.098 & 0.069 & 0.098 & 1.20 & 11.4 \\
Full CMB (2018) + DESI & 0.080 & 0.079 & 0.083 & 1.00 & 9.0 \\
Full CMB (D1) + DESI & 0.068 & 0.059 & 0.072 & 1.48 & 5.0 \\
\bottomrule
\end{tabular}
 \end{table}

These chains are far from converged. The SPT 2018 + DESI chain gives an empirical 95th percentile of $0.084$\,eV in its first half and $0.217$\,eV in its second half, with an effective sample size near five. The full CMB (2018) + DESI chain has split $\hat R\approx1.00$ but an effective sample size near nine. Cobaya's internal multivariate check for that chain reports $R-1\approx22$, and a short independent repeat with another seed fails rank-normalized and tail diagnostics. A single scalar $\hat R$ near one cannot overrule such evidence. Differences between the percentiles in \cref{tab:archived-chains} therefore do not measure the effect of the lens swap. In both archived pairs the D1 value is the lower one, but chains this short cannot establish even that ordering.

\subsection{Multi-chain sampling and the stability of the target}
We then ran four to six independently seeded chains per baseline and lens, with the thresholds of \cref{sec:mnu-inference} fixed in advance. This campaign showed that defining the target, and not only sampling it, required care:
\begin{itemize}[leftmargin=*, itemsep=2pt]
  \item \textbf{Solver version.} Evaluating the same SPT+DESI points with CAMB 1.6.5 and CAMB 2.0.4 changes the total $\chi^2$ by up to about 9. The change varies from point to point by up to about 3.7, and the $\chi^2$ response to a $0.01$\,eV mass step changes by up to about 0.7. A version change therefore alters the shape of the likelihood, not just its normalization, and cannot be treated as a sampling correction.
  \item \textbf{Evaluation-order dependence.} With caching enabled, returning to an earlier CAMB transfer state after evaluating a different cosmology changed the TT spectra slightly. One stored chain row then differed from a fresh evaluation by $0.195$ in $\chi^2$. Forcing a fresh transfer calculation at every point removed the discrepancy in all traced comparisons.
  \item \textbf{Boltzmann robustness and precision.} One CLASS run stopped at a mass of about $0.027$\,eV with a non-finite photon scattering rate at the start of the perturbation integration. We traced this to a finite-difference Jacobian step crossing a helium freeze-out cutoff, and patched it. At tested points, raising CLASS precision settings shifts $\chi^2$ by a few tenths, by different amounts at different points.
\end{itemize}
Before the evaluation-order dependence was found, twelve SPT+DESI chains (six per lens) had passed every pre-declared gate. Their likelihood evaluations used cached transfer states, so those readouts no longer count as qualified. No other cohort of either baseline passed every gate on a fixed numerical target. The unqualified cohort readouts also do not establish the direction of the shift, and we draw no conclusion about it from them.

\subsection{What is and is not supported}
We draw three conclusions.
\begin{enumerate}[leftmargin=*, itemsep=2pt]
  \item We report no value for the change in the $\mnu$ upper limit under the lens swap, neither for SPT+DESI nor for the full baseline. This is a statement about what our sampling supports. It does not show that the limit is unchanged, and it does not contradict \cite{arxiv2601_16277}, whose CMB+DESI and CMB-D1+DESI limits are $\mnu<0.11$\,eV and $\mnu<0.069$\,eV with their own pipeline and SPT lensing included.
  \item A qualified claim of a tightening needs error control for each of the two upper-limit estimates. If each estimated 95th percentile is within $\varepsilon_A$ and $\varepsilon_B$ of its target, the estimated shift is within $\varepsilon_A+\varepsilon_B$ of the true shift. This bound is sharp, and the failure probabilities add without any independence assumption (\cref{prop:readout}). A tightening is then established only if the estimated shift exceeds $\varepsilon_A+\varepsilon_B$. These are bounds on estimation error relative to a fixed likelihood target. Solver-version and numerical-precision effects must be controlled separately.
  \item The small posterior mass near the boundary in our chains is a descriptive property of unconverged chains. At most about $0.22\%$ of retained samples lie below $10^{-3}$\,eV in the four archived chains, and up to about $2\%$ in some later campaign chains. It says nothing either way about a boundary mode in the true posterior. \Cref{sec:boundary-mode} explains why boundary behavior must be read with care.
\end{enumerate}
 \section{Boundary modes under a physical gate}
\label{sec:boundary-mode}

The authors of \cite{arxiv2601_16277} report that with D1 the $\mnu$ posterior is pushed against the prior border at $\mnu=0$. This section asks, in the simplest possible model, when such boundary behavior arises and how it responds to tightening. Every proposition below is proved in Lean~4 with Mathlib. Appendix~\ref{app:formal} maps each statement to its formal counterpart. The model assumes exactly Gaussian, independent constraints and, for posterior statements, a prior that is constant on the allowed region. These are assumptions of the model, and we have not shown that real CMB or BAO likelihoods satisfy them.

\subsection{Two Gaussian constraints and a gate}
Let $m\ge0$ stand for $\mnu$. Two independent constraints contribute Gaussian factors $\exp[-(m-\mu_i)^2/(2\sigma_i^2)]$ with $\sigma_i>0$ and precisions $a=\sigma_1^{-2}$, $b=\sigma_2^{-2}$. Completing the square gives
\begin{equation}
  a(m-\mu_1)^2+b(m-\mu_2)^2 \;=\; (a+b)(m-\mu_\star)^2+\frac{ab}{a+b}(\mu_1-\mu_2)^2,
  \qquad
  \mu_\star=\frac{a\mu_1+b\mu_2}{a+b},\quad \sigma_\star^2=\frac{1}{a+b}.
  \label{eq:toy_combined_gaussian}
\end{equation}

\begin{proposition}[Gated mode]\label{prop:mode}
On $[0,\infty)$ the product of the two factors has the unique maximizer
\begin{equation}
  m_{\mathrm{mode}}=\max(0,\mu_\star).
  \label{eq:toy_trunc_mode}
\end{equation}
The mode is at the boundary exactly when $a\mu_1+b\mu_2\le0$. In particular, if $\mu_1\ge0$ and $\mu_2\ge0$ then $\mu_\star\ge0$. Multiplying both precisions by the same positive constant leaves $\mu_\star$ unchanged.
\end{proposition}

The second statement is the key qualitative point. Tension between two constraints that both prefer positive masses can never move the mode to zero. A boundary mode needs a nonpositive combined mean. Apart from the edge case in which both means are exactly zero, this requires at least one constraint whose effective mean is negative, a pull towards ``negative mass''. The gate turns this pull into a mode at zero, without negative masses ever being physical (\cref{fig:toy}a).

\begin{proposition}[The gated law]\label{prop:law}
Let $\sigma>0$ and $\mu\in\mathbb{R}$. The kernel $\exp[-(m-\mu)^2/(2\sigma^2)]$ restricted to $m\ge0$ has integral $Z=\sigma\sqrt{2\pi}\,\Phi(\mu/\sigma)>0$, where $\Phi$ is the standard normal CDF. The normalized law is the Gaussian $\mathcal N(\mu,\sigma^2)$ restricted to $[0,\infty)$ and divided by $\Phi(\mu/\sigma)$. Its CDF is
\begin{equation}
  F(x)=\frac{\Phi\!\left(\frac{x-\mu}{\sigma}\right)-\Phi\!\left(-\frac{\mu}{\sigma}\right)}{\Phi\!\left(\frac{\mu}{\sigma}\right)}\quad(x\ge0),\qquad F(x)=0\quad(x\le0),
  \label{eq:gated_cdf}
\end{equation}
which is continuous and strictly increasing on $[0,\infty)$. The law gives zero probability to the point $m=0$ and positive probability to every interval $(0,\varepsilon]$ with $\varepsilon>0$. For every $0<q<1$ it has exactly one $q$-quantile, and that quantile is positive.
\end{proposition}

Thus a density mode at zero is not a probability atom at zero, and upper limits remain well defined and strictly positive. Three quantities should be kept apart: the location of the density mode, the posterior mass in a small interval near zero, and the probability of obtaining a boundary mode across repeated experiments.

\subsection{How tightening changes the chance of a boundary mode}
Suppose the second constraint's mean scatters around the first: $\mu_2=\mu_1+\epsilon$ with offset $\epsilon\sim\mathcal N(-d,s^2)$ and $s>0$. A positive $d$ means that the second constraint tends to prefer smaller masses.

\begin{proposition}[Boundary probability]\label{prop:sweep}
The gated mode is at zero exactly when $\epsilon\le-\mu_1\left(1+\sigma_2^2/\sigma_1^2\right)$, so
\begin{equation}
  P(\mathrm{mode}=0)=\Phi\!\left(\frac{d-\mu_1\left(1+\sigma_2^2/\sigma_1^2\right)}{s}\right).
  \label{eq:sweep_probability}
\end{equation}
Hold $\mu_1$, $\sigma_1$, $d$ and $s$ fixed. If $\mu_1>0$, the probability strictly increases as the second constraint tightens ($\sigma_2$ decreases). If $\mu_1<0$, it strictly decreases. If $\mu_1=0$, it does not depend on $\sigma_2$. Tightening both constraints by a common factor leaves the boundary event unchanged.
\end{proposition}

So ``tighter data make boundary modes more common'' holds only in a specific regime: one constraint is tightened, and the other has a positive mean. It is not a general law (\cref{fig:toy}b). The tightening statement also assumes that the offset distribution stays the same while $\sigma_2$ changes. If a new data release changes both the width and the mean of a constraint, as a lens swap may, the proposition does not apply directly.

\begin{figure}[t]
  \centering
  \includegraphics[width=\linewidth]{figures/fig_toy_boundary.pdf}
  \caption{The two-constraint model. (a) With $\mu_1=0.05$, $\sigma_1=0.02$, $\mu_2=-0.02$ and $\sigma_2=0.01$\,eV, the ungated product (dashed) peaks at $\mu_\star=-0.006$\,eV. After the gate, the density (solid) has its mode at $m=0$ but no atom there. (b) Probability of a boundary mode, \cref{eq:sweep_probability}, as the second constraint tightens, for $\mu_1\in\{+0.02,0,-0.02\}$\,eV, $\sigma_1=0.02$\,eV, $d=0$ and $s=0.03$\,eV. Tightening raises the probability when $\mu_1>0$, lowers it when $\mu_1<0$, and leaves it unchanged when $\mu_1=0$.}
  \label{fig:toy}
\end{figure}

\subsection{The finite upper prior and the reading of upper limits}
Posteriors are computed with a finite prior box, here $\mnu\le5$\,eV. The next statement shows what such a box can and cannot hide.

\begin{proposition}[Finite upper gate]\label{prop:upper-gate}
Let $\nu$ be any probability law on $\mathbb{R}$ with CDF $F$ and $F(U)>0$. Conditioning on $m\le U$ gives the CDF $F_U(x)=F(\min(x,U))/F(U)$. For every $x$,
\[
  0\;\le\; F_U(x)-F(x)\;\le\; 1-F(U),
\]
with equality on the right at $x=U$. For the gated Gaussian of \cref{prop:law} and any $U>0$, each level $0<q<1$ has a unique conditioned quantile in $(0,U)$.
\end{proposition}

The error introduced by a finite upper gate is therefore bounded by the probability that the wider posterior assigns above $U$. That probability has to be established separately. A small conditional upper limit does not establish it. As an exact example, take a likelihood equal to $47/10$ on $[0,1/5]$, $3/50$ on $[5,6]$, and zero elsewhere. With a flat prior on $[0,5]$ its 95th percentile is $0.19$\,eV, while with a flat prior on $[0,6]$ it is $31/6\approx5.17$\,eV. We do not suggest that the real likelihood behaves like this. The example only shows why the absence of samples near $5$\,eV is not, by itself, a proof that the prior box is irrelevant.

\begin{proposition}[Comparing two upper limits]\label{prop:readout}
Let $q_A,q_B$ be two upper limits and $\hat q_A,\hat q_B$ their estimates, with $|\hat q_A-q_A|\le\varepsilon_A$ and $|\hat q_B-q_B|\le\varepsilon_B$. Then $|(\hat q_A-\hat q_B)-(q_A-q_B)|\le\varepsilon_A+\varepsilon_B$, and the bound can be attained. If $\hat q_A-\hat q_B>\varepsilon_A+\varepsilon_B$ then $q_A>q_B$. When the estimates are random variables on a common probability space, the probability that the combined bound fails is at most the sum of the two individual failure probabilities, whatever the dependence between the estimates.
\end{proposition}

\subsection{What the model does and does not say about the lens swap}
The model shows that a hard gate can turn a nonpositive combined mean, for example from a moderate negative pull, into a mode at exactly zero, and that whether tightening favors this depends on the sign of the other constraint. It does not show that the SPT lens swap creates such a pull. Nor does it show that the real posteriors are Gaussian or that their constraints are independent. Our own sampled posteriors cannot settle the question (\cref{sec:results-mnu}). The model's practical lesson is about reporting. When the posterior sits against the gate, quantiles of the gated law (\cref{prop:law}) remain meaningful. The location of the mode, in contrast, can switch to or away from zero under modest changes in a constraint's mean.
 \section{Discussion}
\label{sec:discussion}

\subsection{What the audits tell us}
The SPT 2018 and D1 likelihoods do not predict each other well. At the audited fits, parameters chosen by one lens cost the other tens to hundreds of units of $\chi^2$, and the cost depends strongly on direction: testing D1 at a 2018 fit is always the more expensive transfer. Extra D1 multipole coverage, the two releases' different optical-depth priors and default Boltzmann precision change these numbers by a few percent at most. Re-fitting a capped set of nuisance parameters absorbs a large share of the fixed-spectrum disagreement, but not all of it. Where the disagreement sits depends on what else is in the fit. Temperature at $2000\le\ell<3000$ leads both directions with SPT and DESI. Polarization leads the D1-given-2018 direction once the Planck subset and PR4 lensing are added.

Read through the SBT lens, the $\mnu$ shift reported in \cite{arxiv2601_16277} is a change of package under a lens swap, and the two lenses disagree at the level of the packaged likelihoods themselves. Until that disagreement is understood, a shift in a tight limit should be regarded as sensitive to packaging. This conclusion does not depend on the size of the shift, and it does not imply that either release is wrong.

\subsection{Limitations}
\begin{itemize}[leftmargin=*, itemsep=2pt]
  \item \textbf{Optimization.} All audit values are differences between numerically optimized candidates without a global-optimality certificate. Improved source fits moved individual values by up to a few percent. The positivity of $\dchi$ is guaranteed only between the retained candidates.
  \item \textbf{Significance.} The two releases differ in seasons, noise and nuisance models. Without a declared null distribution the $\dchi$ values are comparative diagnostics, not $p$-values. Signed allocations are not independent chi-squares.
  \item \textbf{Nuisance scope.} The SPT+DESI audits and chains fix SPT nuisance parameters at packaged defaults. The full baseline frees only three, and the fixed-spectrum audit re-fits a capped set of twelve. Broader nuisance freedom could change both the audit values and the posteriors.
  \item \textbf{Baseline scope.} Our full baseline omits SPT lensing and differs from the SPT+DESI baseline in Boltzmann code and free foregrounds. Comparisons between the two baselines do not isolate a single cause.
  \item \textbf{Posteriors.} We have no converged $\mnu$ posterior on a qualified numerical target (\cref{sec:results-mnu}). The likelihoods' sensitivity to solver version, caching and precision settings is part of the reason.
  \item \textbf{Released objects only.} The audits operate on released likelihoods. They say where the packaged summaries disagree, not whether calibration, beams, transfer functions, foreground modeling, map-making or covariance construction is responsible.
\end{itemize}

\subsection{Next diagnostics}
The decisive test remains the one proposed in \cite{arxiv2601_16277}: reprocess the 2018 raw data through the D1 pipeline. Short of that, three steps follow from this work.
\begin{enumerate}[leftmargin=*, itemsep=2pt]
  \item \textbf{Predictive audits.} Replace best-fit transfer with posterior predictive checks across lenses \cite{gelman_meng_stern1996_ppc}, with a declared null that accounts for differing noise and seasons. This would turn the localization maps into calibrated statements.
  \item \textbf{Controlled targets.} Fix and record the solver version, precision settings and caching behavior. Verify stored chain rows against fresh evaluations. Only then sample to the quantile precision required by \cref{prop:readout}, given the size of the shift to be resolved.
  \item \textbf{Targeted rewrites.} Where the mismatch localizes (EE at $2000\le\ell<3000$ in the full baseline, TT in the same range with SPT and DESI), test small template or calibration corrections fitted on one lens and evaluated on the other. A correction selected to fit the residual it removes improves the fit by construction. Its value must be judged on held-out data.
\end{enumerate}

\subsection{Practical recommendation}
When a likelihood update moves a tight bound, we suggest reporting, alongside the bound: the directional mismatch in both directions; its signed, covariance-respecting localization by spectrum and multipole, for each baseline used; at least one support control; and posterior diagnostics with quantile-precision targets matched to the size of the claimed shift. These are inexpensive compared with the inference itself, and they change how a moving bound should be read.
 \section{Conclusion}
\label{sec:conclusion}

We examined the reported tightening of the $\mnu$ limit when SPT-3G 2018 is replaced by SPT-3G D1 by asking first whether the two releases agree with each other. Treating each release as a lens in the Six Birds sense, we measured directional mismatches at fitted candidates and split their SPT quadratic part exactly over spectra and multipoles. The releases disagree strongly and asymmetrically. The location of the disagreement depends on the baseline and the direction of transfer: high-$\ell$ temperature leads with SPT and DESI, and high-$\ell$ polarization leads the D1-given-2018 direction once a Planck subset is added. Multipole-support, optical-depth-prior and precision controls change the totals by only a few percent. Our $\mnu$ chains do not meet pre-declared convergence and precision criteria, and the likelihoods proved sensitive to solver details, so we report no magnitude for the tightening. A machine-checked two-constraint model shows that a boundary mode at $\mnu=0$ requires a nonpositive combined mean, which tension between constraints that both prefer positive masses cannot produce, and that tightening one constraint can make such a mode either more or less likely. Our central message is methodological: a bound that moves under a lens swap should be reported together with directional audits, their localization and controls, and posterior diagnostics strong enough to resolve the shift being claimed.
 
\appendix
\section{Formal verification}
\label{app:formal}

The mathematical statements of \cref{sec:sbt-framework,sec:methods,sec:boundary-mode} are proved in the Lean~4 library \path{lean/trunc_gauss_proof} (Lean 4.27.0 with a pinned Mathlib revision \cite{lean4,mathlib}). The library exports 97 public declarations. An automated audit (\path{AuditAxioms.lean}) prints the transitive axiom dependencies of each one. All of them depend only on Lean's standard axioms \texttt{propext}, \texttt{Classical.choice} and \texttt{Quot.sound}, with no \texttt{sorry}, \texttt{admit} or added axioms. Running \texttt{make math\_check} from the repository root builds the library, runs the axiom audit and runs the numerical regression tests. \Cref{tab:lean} maps each statement in the text to its formal counterpart.

\begin{table}[h]
  \centering
  \caption{Statements in the text and their Lean counterparts (module and representative declarations).}
  \label{tab:lean}
  \footnotesize
  \setlength{\tabcolsep}{4pt}
  \begin{tabular}{@{}>{\raggedright\arraybackslash}p{0.27\linewidth} >{\raggedright\arraybackslash}p{0.21\linewidth} >{\raggedright\arraybackslash}p{0.46\linewidth}@{}}
    \toprule
    Statement & Module & Representative declarations \\
    \midrule
    Sign and sum rule of $\dchi$ (\cref{sec:sbt-framework}) & \texttt{Audit} & \texttt{directionalDelta\_\allowbreak{}nonneg}, \texttt{directionalDelta\_\allowbreak{}sum} \\
    Signed covariance ledger, \cref{eq:signed_alloc,eq:block_dchi} & \texttt{CovarianceLedger} & \texttt{covariance\_\allowbreak{}solve\_\allowbreak{}eq\_\allowbreak{}inverse}, \texttt{covariance\_\allowbreak{}solve\_\allowbreak{}allocation\_\allowbreak{}sum}, \texttt{groupedQuadraticDifference\_\allowbreak{}with\_\allowbreak{}remainder}, \texttt{residualQuadratic\_\allowbreak{}inverse\_\allowbreak{}nonneg} \\
    Completed square and gated mode, \cref{prop:mode} & \texttt{TruncGauss}, \texttt{Mode}, \texttt{Product} & \texttt{combine\_\allowbreak{}completed\_\allowbreak{}square}, \texttt{gaussianProduct\_\allowbreak{}unique\_\allowbreak{}mode}, \texttt{gaussianProduct\_\allowbreak{}boundary\_\allowbreak{}max\_\allowbreak{}iff}, \texttt{combinedMean\_\allowbreak{}nonneg}, \texttt{combinedMean\_\allowbreak{}nonpos\_\allowbreak{}iff}, \texttt{combinedMean\_\allowbreak{}common\_\allowbreak{}scale} \\
    Normalizer, law, CDF and quantiles, \cref{prop:law} & \texttt{Normalization}, \texttt{Gaussian\-Normalization}, \texttt{GaussianMeasure}, \texttt{GaussianGatedCDF}, \texttt{GaussianQuantile} & \texttt{truncatedNormalizer\_\allowbreak{}eq\_\allowbreak{}cdf}, \texttt{truncatedGaussianLaw\_\allowbreak{}isProbability}, \texttt{truncatedGaussianLaw\_\allowbreak{}singleton}, \texttt{truncatedGaussianLaw\_\allowbreak{}boundary\_\allowbreak{}interval\_\allowbreak{}pos}, \texttt{truncatedGaussianLaw\_\allowbreak{}cdf\_\allowbreak{}of\_\allowbreak{}nonneg}, \texttt{truncatedGaussianLaw\_\allowbreak{}cdf\_\allowbreak{}strictMonoOn}, \texttt{truncatedGaussianLaw\_\allowbreak{}existsUnique\_\allowbreak{}quantile} \\
    Boundary probability and its monotonicity, \cref{prop:sweep} & \texttt{Tightening}, \texttt{GaussianSweep} & \texttt{product\_\allowbreak{}boundary\_\allowbreak{}iff\_\allowbreak{}offset\_\allowbreak{}event}, \texttt{offsetBoundaryEvent\_\allowbreak{}common\_\allowbreak{}scale}, \texttt{standardDeviation\-Boundary\-Gaussian\-Probability\_\allowbreak{}cdf}, \texttt{\dots\_\allowbreak{}strictAnti}, \texttt{\dots\_\allowbreak{}strictMono}, \texttt{\dots\_\allowbreak{}zero} \\
    Finite upper gate, \cref{prop:upper-gate} & \texttt{UpperGate} & \texttt{upperConditionedLaw\_\allowbreak{}cdf}, \texttt{upperConditionedLaw\_\allowbreak{}cdf\_\allowbreak{}error}, \texttt{upperConditionedLaw\_\allowbreak{}cdf\_\allowbreak{}error\_\allowbreak{}at\_\allowbreak{}gate}, \texttt{upperConditionedGaussian\_\allowbreak{}existsUnique\_\allowbreak{}quantile} \\
    Comparing two upper limits, \cref{prop:readout} & \texttt{Readout} & \texttt{tighteningReadout\_\allowbreak{}error\_\allowbreak{}le}, \texttt{tighteningReadout\_\allowbreak{}pos\_\allowbreak{}of\_\allowbreak{}error\_\allowbreak{}bound}, \texttt{tighteningReadout\_\allowbreak{}error\_\allowbreak{}bound\_\allowbreak{}sharp}, \texttt{tighteningReadout\_\allowbreak{}failure\_\allowbreak{}risk\_\allowbreak{}le} \\
    \bottomrule
  \end{tabular}
\end{table}

\paragraph{Scope of the formal results.}
The proofs concern exact real numbers and exact probability measures. They take as hypotheses the conditions stated in the text: Gaussian factors, positive widths, independence of the two constraints, a constant prior on the gate when a product is read as a posterior, a fixed offset law in \cref{prop:sweep}, and, for the audit sign, the stated likelihood ordering. They do not certify the floating-point implementation, show that the cosmological likelihoods are Gaussian, prove MCMC convergence, or establish that any optimizer has found a global maximum. The exact-rational counterexample after \cref{prop:upper-gate} was checked separately and is not part of the Lean library.
 \section{Reproducibility and numerical safeguards}
\label{app:reproducibility}

\paragraph{Code and artifacts.}
The analysis code, run configurations, Lean library and manuscript sources are available at \url{https://github.com/ioannist/six-birds-neutrino}. The audit figures and the two numerical tables (\cref{tab:audit-summary,tab:archived-chains}) are generated directly from stored run bundles by \path{scripts/make_paper_figures.py}. Its output, \path{paper/figures/figure_sources.json}, records the plotted totals with their source bundles and the parameters of the analytic illustration. The two-constraint figure is evaluated afresh from the analytic formulas, using the same numerical routines as the toy-model runs and the parameters stated in its caption. Localization values and percentages quoted in the text come from the same bundles and from the multipole-range check listed below. The baseline table summarizes \cite{arxiv2601_16277}. The main bundles, all under \path{runs/}, are:
\begin{itemize}[leftmargin=*, itemsep=1pt]
  \item cosmological audits: \texttt{20261003\_\allowbreak{}math\_\allowbreak{}review\_\allowbreak{}cosmological\_\allowbreak{}spt\_\allowbreak{}desi}, \texttt{\dots\_\allowbreak{}spt\_\allowbreak{}desi\_\allowbreak{}staging}, \texttt{\dots\_\allowbreak{}cmb\_\allowbreak{}desi}, \texttt{\dots\_\allowbreak{}cmb\_\allowbreak{}desi\_\allowbreak{}warm}, \texttt{\dots\_\allowbreak{}cmb\_\allowbreak{}desi\_\allowbreak{}staging}, \texttt{\dots\_\allowbreak{}cmb\_\allowbreak{}desi\_\allowbreak{}tau\_\allowbreak{}control}, \texttt{\dots\_\allowbreak{}cmb\_\allowbreak{}desi\_\allowbreak{}accuracy} and \texttt{\dots\_\allowbreak{}cmb\_\allowbreak{}desi\_\allowbreak{}accuracy\_\allowbreak{}improved\_\allowbreak{}source};
  \item fixed-spectrum audits: \texttt{20261003\_\allowbreak{}math\_\allowbreak{}review\_\allowbreak{}profiled\_\allowbreak{}multistart}, \texttt{\dots\_\allowbreak{}profiled\_\allowbreak{}staging\_\allowbreak{}multistart} and \texttt{\dots\_\allowbreak{}verified\_\allowbreak{}native\_\allowbreak{}covariance};
  \item multipole-range totals: \texttt{20261003\_\allowbreak{}math\_\allowbreak{}review\_\allowbreak{}localization\_\allowbreak{}range\_\allowbreak{}check};
  \item archived chain diagnostics: \texttt{20261003\_\allowbreak{}math\_\allowbreak{}review\_\allowbreak{}chain\_\allowbreak{}diagnostics};
  \item two-constraint model (numerical checks of the routines used for \cref{fig:toy}): \texttt{20261003\_\allowbreak{}math\_\allowbreak{}review\_\allowbreak{}toy\_\allowbreak{}cancellation}.
\end{itemize}
Each cosmological bundle includes an independent re-evaluation of every endpoint with a freshly built model (\texttt{accounting\_\allowbreak{}verification.json}). A complete record of the checks behind this paper, including the multi-chain sampling campaign and its solver controls, is in \path{docs/findings/math_review.md}.

\paragraph{Capped nuisance sets in the fixed-spectrum audit.}
When D1 is the test lens, the re-fitted set is \texttt{Tcal\_\allowbreak{}ext150}, \texttt{Tcal\_\allowbreak{}rel90}, \texttt{Tcal\_\allowbreak{}rel220}, \texttt{Ecal\_\allowbreak{}ext150}, \texttt{Ecal\_\allowbreak{}rel90}, \texttt{Ecal\_\allowbreak{}rel220}, \texttt{TT\_\allowbreak{}CIBClustering\_Alpha}, \texttt{TT\_\allowbreak{}CIB\_150x150}, \texttt{TT\_\allowbreak{}CIB\_150x220}, \texttt{TT\_\allowbreak{}CIB\_220x220}, \texttt{TT\_\allowbreak{}GalCirrus\_Alpha} and \texttt{TT\_\allowbreak{}GalCirrus\_Amp}. When 2018 is the test lens, it is \texttt{TT\_\allowbreak{}tSZ\_CIB\_Corr\_Amp}, \texttt{TT\_\allowbreak{}tSZ\_CIB\_corr}, \texttt{Ecal150}, \texttt{Ecal220}, \texttt{Ecal90}, \texttt{TT\_\allowbreak{}CIBClustering\_Alpha}, \texttt{TT\_\allowbreak{}CIBClustering\_Amp}, \texttt{TT\_\allowbreak{}CIBClustering\_Beta}, \texttt{TT\_\allowbreak{}GalCirrus\_Alpha}, \texttt{TT\_\allowbreak{}GalCirrus\_Amp}, \texttt{TT\_\allowbreak{}GalCirrus\_Beta} and \texttt{TT\_\allowbreak{}Poisson\_150x150}. Both sets are capped at twelve parameters to keep the run time tractable. The shared parameters are the tSZ and kSZ amplitudes.

\paragraph{Numerical safeguards.}
The pipeline enforces the following properties, each covered by regression tests (245 Python tests in total):
\begin{itemize}[leftmargin=*, itemsep=1pt]
  \item \emph{Covariances} must be finite, symmetric and positive definite. Only round-off asymmetry is symmetrized, and singular matrices are rejected rather than pseudo-inverted. Solves, allocations and their sums that cannot be represented in floating point are refused rather than returned as infinities or NaNs.
  \item \emph{Quantiles} use exact integer or rational cumulative mass, so they do not depend on how samples are grouped into rows.
  \item \emph{Chains} are combined only when their likelihoods, priors, sampler temperature (which must be one), Boltzmann code version and numerical settings all match. Column names in chain files must be unique and unambiguous.
  \item \emph{Two-constraint model}: the combination of \cref{eq:toy_combined_gaussian}, the boundary probability of \cref{eq:sweep_probability} and the gated density remain accurate under extreme weights, near-cancelling means and far tails. They fall back to exact rational arithmetic, or to a bounded asymptotic tail expansion, when floating point would lose the answer.
  \item \emph{Toy template rewrite}: a covariance-weighted projection of a template out of a residual is computed in whitened coordinates, with exact accumulation. When the template is chosen from the residual itself, the improvement in fit is guaranteed by construction and is reported as descriptive only.
\end{itemize}
These safeguards are tested numerical properties, not floating-point error certificates.
 \section{Version history}
\label{app:versions}

\paragraph{\manuscriptversion\ (\manuscriptdate).}
This version follows a complete mathematical, numerical and formal review of the analysis. The main changes from Version 1 are:
\begin{itemize}[leftmargin=*, itemsep=1pt]
  \item The Version 1 tightening magnitudes (about $0.109$\,eV with SPT and DESI, and about $0.012$\,eV with the Planck-subset baseline) were derived from single, unconverged chains. They are withdrawn as unsupported. They have not been shown to be false. No replacement magnitude is given.
  \item Localization previously inverted each block of the covariance separately. It now uses the exact signed allocation of \cref{sec:ledger}. The audits were re-run, including cosmological audits with a shared Boltzmann calculation and multi-start fits. The claim that high-$\ell$ TT dominates the mismatch in general is replaced by the baseline- and direction-dependent statement of \cref{sec:results-audits}.
  \item The fixed-spectrum profiled audit now releases the shared parameters in the reference fit and uses multiple starting points. Its values change from $859.6$/$369.5$ to $849.6$/$414.4$.
  \item The boundary-mode discussion now states its hypotheses explicitly. It shows that a boundary mode requires a nonpositive combined mean, which two constraints with positive means cannot produce, corrects the claim that tightening generally increases boundary-mode frequency (\cref{prop:sweep}), and adds the finite-upper-gate and readout results. The Lean library grew from two lemmas to 97 audited declarations.
  \item Chain summaries now use the inverse weighted empirical CDF with exact holding-time weights, together with chronological split diagnostics. For the same stored samples, this changes the archived SPT D1 + DESI 95th percentile from $0.085$ to $0.098$\,eV, and it shows that none of the archived chains is converged (\cref{tab:archived-chains}).
  \item The abstract and main text were rewritten. The audit figures and tables are now generated directly from the verified run bundles, and the two-constraint figure is re-evaluated from the analytic formulas.
\end{itemize}

\paragraph{Version 1 (2 March 2026).}
First public version, posted on Research Square \cite{tsiokos2026_neutrino_v1}.
 
\bibliographystyle{unsrt}
\bibliography{references}

\end{document}
